\documentclass[11pt,a4paper]{article}
\usepackage{fontspec,amsmath,amssymb,amsthm,microtype}
\setmainfont{lmroman10-regular.otf}[BoldFont=lmroman10-bold.otf,ItalicFont=lmroman10-italic.otf,BoldItalicFont=lmroman10-bolditalic.otf]
\usepackage[margin=25mm]{geometry}
\usepackage[unicode,hidelinks]{hyperref}
\setlength{\emergencystretch}{3em}
\setlength{\parskip}{0.4em}
\setlength{\parindent}{0pt}
\allowdisplaybreaks[2]
\raggedbottom
\title{\texorpdfstring{Construct the averaged Hessian on a segment}{Construct the averaged Hessian on a segment}}
\author{Omer Deniz Akyildiz\\[0.4em]\normalsize Prepared with AI assistance}
\date{}
\begin{document}\maketitle
Draft mathematical exposition; these notes may contain mistakes.


The mean Hessian is a Bochner integral of the continuous derivative of the force along the segment between two positions. The module derives its self-adjointness, lower and upper quadratic-form bounds and exact representation of the force difference. Thus the operator used in the spectral contraction proof is a concrete object, not an assumed matrix already chosen to make the desired equation hold.



The averaged operator is not an arbitrary matrix satisfying spectral bounds. It is constructed from the Hessian along the displacement segment and is proved to reproduce the exact force difference.


\section{\texorpdfstring{Define the segment average and commute it with a pairing}{Define the segment average and commute it with a pairing}}


Set \(\bar H(x,y)=\int_0^1H(y+t(x-y))dt\). Continuity of \(H\) makes the operator-valued integrand continuous and integrable on the compact segment. A continuous linear map commutes with a Bochner integral, so applying the operator to \(w\) and pairing with \(z\) gives \(\langle z,\bar H(x,y)w\rangle=\int_0^1\langle z,H(y+t(x-y))w\rangle dt\). First, the argument commutes operator evaluation through the interval integral, then commutes the continuous linear inner-product functional. The operator type is complete because the underlying Hilbert space is complete.



The average is taken along the straight line joining the two positions. Pairing with fixed vectors turns the operator-valued integral into an ordinary scalar integral. This permits both symmetry and quadratic-form bounds to be through scalar identities on the segment.


\section{\texorpdfstring{Average symmetry and both spectral inequalities}{Average symmetry and both spectral inequalities}}


If every \(H(p)\) is self-adjoint, the pairing identity moves both sides into scalar integrals, where their integrands agree pointwise. If \(m\|z\|^2\le\langle z,H(p)z\rangle\le M\|z\|^2\), integral monotonicity bounds the average between the integrals of these two constants. The interval has length one, leaving exactly the original constants. The argument supplies continuity and interval-integrability for each quadratic-form integrand and the constant comparison function, applies monotonicity of integration, and simplifies the constant integral. These results provide the symmetry and spectral premises needed by the transformed-energy proof.



For every vector \(u\), integrate \(m\|u\|^2\le\langle H_su,u\rangle\le M\|u\|^2\) over an interval of length one. The bounds remain \(m\) and \(M\). Symmetry is preserved because the equality of the two mixed pairings holds at each segment point.


\section{\texorpdfstring{The fundamental theorem identifies the force difference}{The fundamental theorem identifies the force difference}}


Assume \(Dg(p)=H(p)\). The affine path \(\ell(t)=y+t(x-y)\) has derivative \(x-y\), so the chain rule gives \((g\circ\ell)'(t)=H(\ell(t))(x-y)\). Integrate this derivative from zero to one. The fundamental theorem yields \(\int_0^1H(\ell(t))(x-y)dt=g(x)-g(y)\). Commuting the operator evaluation with the average integral gives \(\bar H(x,y)(x-y)=g(x)-g(y)\). The argument constructs the affine-path derivative, composes the force derivative at each segment point, verifies integrability from continuity, and simplifies the endpoints. No coordinatewise mean-value theorem or unjustified single Hessian point is used.



Apply the fundamental theorem of calculus to \(s\mapsto\nabla f(y+s(x-y))\). Its derivative is \(\nabla^2f(y+s(x-y))(x-y)\). Integrating from zero to one gives the force difference, proving that the averaged operator is the coefficient in the synchronous difference equation.


\end{document}
